\section{Background and Related Work}
\label{sec:background}

\subsection{NS-3 Simulator Architecture}

NS-3 is a discrete-event network simulator designed for research and educational purposes~\cite{ns3}. Unlike its predecessor NS-2, NS-3 is implemented entirely in C++ with optional Python bindings, providing strong typing, better performance, and closer alignment with real-world protocol implementations.

\subsubsection{Discrete Event Simulation Model}

NS-3 employs a discrete-event simulation paradigm where the system state changes only at discrete points in time when events occur. The core scheduling mechanism maintains a future event list (FEL) ordered by simulation time. The simulator processes events in chronological order, ensuring strict causality: an event at time $t$ cannot affect events at time $t' < t$.

The simulation proceeds as follows:
\begin{enumerate}
    \item Extract the earliest event $e$ from the FEL
    \item Advance simulation time to $e$.time
    \item Execute $e$.handler(), which may:
    \begin{itemize}
        \item Modify simulation state
        \item Schedule new events in the FEL
        \item Generate output/statistics
    \end{itemize}
    \item Repeat until FEL is empty or stop time is reached
\end{enumerate}

This model provides reproducible, deterministic results independent of real-time execution speed, making it ideal for controlled experiments.

\subsubsection{WiFi Module Architecture}

NS-3's WiFi module implements the IEEE 802.11 protocol stack with high fidelity to the standard. The architecture follows a layered design, as illustrated in Figure~\ref{fig:ns3-wifi-stack}:

\begin{figure}[htbp]
\centering
\begin{verbatim}
┌─────────────────────────────────────┐
│     Application Layer               │
│  (UDP/TCP Applications)             │
└─────────────────┬───────────────────┘
                  │
┌─────────────────▼───────────────────┐
│     Network Layer (IPv4/IPv6)       │
└─────────────────┬───────────────────┘
                  │
┌─────────────────▼───────────────────┐
│     WiFi MAC (802.11 DCF/EDCA)     │
│  - Frame aggregation (A-MPDU)       │
│  - Block ACK                        │
│  - Rate adaptation                  │
└─────────────────┬───────────────────┘
                  │
┌─────────────────▼───────────────────┐
│     WiFi PHY Layer                  │
│  - OFDM/DSSS modulation            │
│  - Error rate models                │
│  - Energy detection                 │
└─────────────────┬───────────────────┘
                  │
┌─────────────────▼───────────────────┐
│     WiFi Channel                    │
│  - Propagation loss models          │
│  - Delay models                     │
│  - Multi-device broadcast           │
└─────────────────────────────────────┘
\end{verbatim}
\caption{NS-3 WiFi protocol stack architecture}
\label{fig:ns3-wifi-stack}
\end{figure}

\paragraph{YansWifiChannel}
The \texttt{YansWifiChannel} class (Yet Another Network Simulator WiFi Channel) is NS-3's primary WiFi channel implementation. When a device transmits:

\begin{lstlisting}[style=cpp, caption={YansWifiChannel transmission logic (simplified)}]
void YansWifiChannel::Send(Ptr<YansWifiPhy> sender,
                           Ptr<const WifiPpdu> ppdu,
                           double txPowerDbm)
{
    // Get sender position
    Ptr<MobilityModel> senderMobility = sender->GetMobility();
    
    // Iterate over all PHYs on this channel
    for (auto receiver : m_phyList)
    {
        if (receiver == sender) continue; // Skip self
        
        // Calculate propagation delay
        Ptr<MobilityModel> receiverMobility = 
            receiver->GetMobility();
        Time delay = m_delay->GetDelay(
            senderMobility, receiverMobility);
        
        // Calculate received power
        double rxPowerDbm = m_loss->CalcRxPower(
            txPowerDbm, senderMobility, receiverMobility);
        
        // Schedule reception at receiver
        Simulator::Schedule(delay, 
            &YansWifiPhy::StartReceivePreamble,
            receiver, ppdu, rxPowerDbm);
    }
}
\end{lstlisting}

This broadcast nature—where each transmission invokes calculations for every potential receiver—creates the O(n²) complexity that limits scalability.

\subsection{MPI and Distributed Simulation}

\subsubsection{Message Passing Interface (MPI)}

MPI is a standardized message-passing library interface specification, implemented by various libraries (OpenMPI, MPICH, Intel MPI)~\cite{mpi-standard}. MPI provides:

\begin{itemize}
    \item \textbf{Point-to-point communication}: \texttt{MPI\_Send()}, \texttt{MPI\_Recv()}
    \item \textbf{Collective operations}: \texttt{MPI\_Broadcast()}, \texttt{MPI\_Reduce()}
    \item \textbf{Synchronization primitives}: \texttt{MPI\_Barrier()}
    \item \textbf{Process topology}: Rank identification, communicator groups
\end{itemize}

MPI programs execute as multiple processes (ranks) that communicate via message passing. Each rank has a unique identifier from 0 to $n-1$ where $n$ is the total process count.

\subsubsection{Parallel Discrete Event Simulation (PDES)}

Distributing discrete event simulation across multiple machines requires solving the time synchronization problem. Two main approaches exist:

\paragraph{Conservative Synchronization}
Conservative protocols~\cite{chandy-misra} guarantee causality by processing events only when it's safe—i.e., when no earlier event could arrive from other processors. This requires "lookahead": knowledge of the minimum timestamp of future messages. The Chandy-Misra-Bryant algorithm uses null messages to communicate time advancement.

\paragraph{Optimistic Synchronization}
Optimistic protocols like Time Warp~\cite{jefferson-timewarp} allow processors to execute events speculatively. If a causality violation is detected (an event arrives with an earlier timestamp than already-processed events), the simulator rolls back state and re-executes. This requires state saving and rollback mechanisms but can achieve better performance when violations are rare.

\subsubsection{NS-3 MPI Support}

NS-3 provides distributed simulation capability through its \texttt{MpiInterface} class and \texttt{DistributedSimulator}~\cite{ns3-mpi}. The implementation uses:

\begin{itemize}
    \item \textbf{Conservative synchronization}: Null message protocol with configurable lookahead
    \item \textbf{Logical processes}: Each MPI rank runs an independent \texttt{DistributedSimulator} instance
    \item \textbf{Time synchronization}: LBTS (Lower Bound on Time Stamp) calculation ensures causality
    \item \textbf{Remote channels}: \texttt{PointToPointRemoteChannel} for distributed wired links
\end{itemize}

Key NS-3 MPI concepts:

\begin{lstlisting}[style=cpp, caption={NS-3 MPI initialization example}]
#include "ns3/mpi-interface.h"

int main(int argc, char *argv[])
{
    // Initialize MPI
    MpiInterface::Enable(&argc, &argv);
    
    uint32_t systemId = MpiInterface::GetSystemId();
    uint32_t systemCount = MpiInterface::GetSize();
    
    // Create topology specific to this rank
    if (systemId == 0) {
        // Rank 0 logic
    } else {
        // Other ranks logic
    }
    
    // Simulation runs with time synchronization
    Simulator::Stop(Seconds(simTime));
    Simulator::Run();
    Simulator::Destroy();
    
    // Clean up MPI
    MpiInterface::Disable();
    return 0;
}
\end{lstlisting}

\subsection{Related Work}

\subsubsection{Distributed Wired Network Simulation}

NS-3's \texttt{PointToPointRemoteChannel} enables distributed simulation of wired networks~\cite{ns3-mpi}. This implementation partitions topology across ranks, with remote channels connecting nodes on different machines. Point-to-point links have well-defined endpoints, making partitioning straightforward. Our WiFi implementation draws inspiration from this design but must address WiFi's broadcast nature.

\subsubsection{Parallel Wireless Simulation}

Several systems have addressed parallel wireless simulation:

\paragraph{JiST/SWANS}
Java in Simulation Time (JiST) with Scalable Wireless Ad hoc Network Simulator (SWANS)~\cite{barr-jist} achieved high scalability through:
\begin{itemize}
    \item Rewriting Java bytecode to intercept blocking operations
    \item Hierarchical spatial indexing for efficient neighbor discovery  
    \item Field partitioning to distribute devices spatially
\end{itemize}

However, JiST's approach requires Java and doesn't integrate with NS-3's C++ codebase.

\paragraph{GTNetS}
Georgia Tech Network Simulator~\cite{gtnet-parallel} explored optimistic synchronization for wireless networks. It demonstrated speedups but required careful tuning to avoid excessive rollbacks in dense wireless scenarios where many events occur in tight temporal proximity.

\paragraph{PRIME SSFNet}
PRIME's SSFNet~\cite{prime-wireless} used conservative synchronization with spatial partitioning. Wireless transmissions crossing partition boundaries required inter-processor communication. Performance degraded when mobile devices frequently crossed boundaries.

\subsubsection{NS-3 LTE MPI Implementations}

The NS-3 LTE module includes experimental MPI support~\cite{ns3-lte-mpi}. LTE's cellular architecture naturally partitions: eNodeBs (base stations) can be distributed across ranks with minimal inter-rank communication for handovers. This provides a reference for integrating MPI with NS-3 wireless modules but doesn't directly address WiFi's peer-to-peer broadcast model.

\subsection{Gap Analysis}

Existing approaches have limitations for large-scale WiFi simulation:

\begin{itemize}
    \item \textbf{Spatial partitioning} works poorly for dense WiFi deployments where all devices are within communication range (e.g., vehicular networks, enterprise WiFi)
    \item \textbf{Optimistic synchronization} suffers from frequent rollbacks due to WiFi's rapid event rates and broadcast nature
    \item \textbf{Existing NS-3 MPI support} targets wired networks; no production-ready WiFi MPI implementation exists
    \item \textbf{Observability gap}: No prior work addresses PCAP capture in distributed wireless simulation
\end{itemize}

Our work fills these gaps by embracing WiFi's centralized broadcast nature, using conservative synchronization aligned with NS-3's existing MPI framework, and introducing a hybrid approach for observability.
